%% kapa: exam=1 round=20 year=2009 period=2 form=A subject=appraisal_law subject_ko=부동산관계법규 seq=17 number=57
%% source: Q-Net 감정평가사 기출문제 (cst00309) artlSeq=5024178 file=extracted/1008/5024178/2009년 제20회 감정평가사 1차시험 기출문제/2교시.hwp
%% conversion: kordoc 4.24.51 -> tools/gen1.py (auto); review=proofread-sonnet; flags=none
\begin{question}{57}
부동산 가격공시 및 감정평가에 관한 법령상 표준지공시지가에 관한 설명으로 옳지 \uline{않은} 것은?
\begin{choices}
  \item 토지시장의 지가정보를 제공한다.
  \item 국가ㆍ지방자치단체 등의 기관이 그 업무와 관련하여 지가를 산정할 때 기준이 된다.
  \item 일반적인 토지거래의 지표가 된다.
  \item 판례는 개별공시지가결정과는 달리 표준지공시지가결정에는 처분성이 없다고 한다.
  \item 감정평가업자가 개별적으로 토지를 감정평가할 때 그 기준이 된다.
\end{choices}
\end{question}
\answer{4}
